% Меняйте метаданные здесь. Для разных работ используйте отдельные копии проекта.
\newcommand{\paperUniversity}{%
  «Московский государственный технический университет\\
  имени Н.\,Э.~Баумана\\
  (национальный исследовательский университет)»\\
  (МГТУ им. Н.\,Э.~Баумана)}
\newcommand{\paperMinistry}{Министерство науки и высшего образования Российской Федерации}
\newcommand{\paperFaculty}{ИУ «ИНФОРМАТИКА И СИСТЕМЫ УПРАВЛЕНИЯ»}
\newcommand{\paperDepartment}{ИУ7 «ПРОГРАММНОЕ ОБЕСПЕЧЕНИЕ ЭВМ И ИНФОРМАЦИОННЫЕ ТЕХНОЛОГИИ»}
\newcommand{\paperTitle}{Тема работы}
\newcommand{\paperAuthor}{И.\,О. Фамилия}
\newcommand{\paperGroup}{Группа}
\newcommand{\paperSupervisor}{И.\,О. Фамилия}
\newcommand{\paperConsultant}{}
\newcommand{\paperEmblem}{img/bmstu-emblem.jpg}
\newcommand{\paperYear}{\number\year}

% Путь относительно docs/ к PDF с утверждённым титулом и листами задания.
% Пустое значение включает редактируемый титул.
\newcommand{\paperFrontmatterPDF}{}
\newcommand{\paperFrontmatterPages}{-}
% Пустое значение продолжает счёт страниц. Для режима образца задайте 4.
\newcommand{\paperFirstPage}{}
\newif\ifpaperAbbreviations
\newif\ifpaperAppendix
\providecommand{\paperProfileDefaults}{}
\paperProfileDefaults
